\documentclass[12pt,a4paper]{extarticle}
\usepackage[margin=1.0cm]{geometry}
\usepackage[T1]{fontenc}
\usepackage{amsmath,amssymb}
\usepackage{multicol}
\usepackage{xcolor}
\usepackage[most]{tcolorbox}
\usepackage{enumitem}
\setlength{\parskip}{0.30em}
\setlength{\parindent}{0pt}
\pagestyle{empty}

\definecolor{cK}{HTML}{1B5E20}\definecolor{cKb}{HTML}{E8F5E9}
\definecolor{cT}{HTML}{B71C1C}\definecolor{cTb}{HTML}{FFEBEE}
\definecolor{cQ}{HTML}{1F4E79}\definecolor{cQb}{HTML}{F0F5FA}

\newtcolorbox{kb}[1][]{enhanced,colback=cKb,colframe=cK,boxrule=0.6pt,
  left=3pt,right=3pt,top=2pt,bottom=2pt,arc=1pt,
  fonttitle=\bfseries\footnotesize,coltitle=black,title={#1}}
\newtcolorbox{tb}[1][]{enhanced,colback=cTb,colframe=cT,boxrule=0.6pt,
  left=3pt,right=3pt,top=2pt,bottom=2pt,arc=1pt,
  fonttitle=\bfseries\footnotesize,title={#1}}
\newtcolorbox{qb}[1][]{enhanced,colback=cQb,colframe=cQ,boxrule=0.6pt,
  left=3pt,right=3pt,top=2pt,bottom=2pt,arc=1pt,
  fonttitle=\bfseries\footnotesize,coltitle=black,title={#1}}

\newcommand{\dd}{\mathrm{d}}
\newcommand{\half}{\tfrac{1}{2}}
\newcommand{\si}{\sigma_{\mathrm{imp}}}

\begin{document}
\begin{center}
{\large\bfseries MAC F244 Viva: what every member must be able to say}\\
{\footnotesize Volatility as a quadratic variation, NIFTY 50. Read this twice before you walk in.}
\end{center}
\vspace{-0.3em}

\begin{multicols}{2}

\begin{kb}[The project in three sentences]
\footnotesize
Black--Scholes needs a volatility nobody can observe. It\^o's lemma applied to $\log S$ shows
that volatility \emph{is} the quadratic variation of the log price, so it can be measured from
one path instead of guessed. We verify that measurement where the answer is known, apply it to
eleven years of NIFTY~50, show the options market over-forecasts volatility, and confirm that
the resulting profit is exactly the one It\^o predicts.
\end{kb}

\textbf{\footnotesize The identity everything rests on}
\vspace{-0.4em}
\[[\log S]_T=\int_0^T\sigma_t^2\,\dd t,\qquad RV_n=\sum(\Delta\log S)^2\]
\vspace{-1.2em}
\textbf{\footnotesize The P\&L identity}
\vspace{0.5em}
\[\Pi_T=\half\int_0^T e^{r(T-t)}S_t^2\Gamma_t\big(\si^2-\sigma_t^2\big)\dd t\]
\vspace{-1.0em}

{\footnotesize Derived from self-financing $+$ It\^o $+$ the Black--Scholes PDE. The $\dd W$
terms cancel because $\Delta=V_S$.}

\begin{kb}[The numbers. Know these cold.]
\scriptsize
\begin{tabular}{@{}ll@{}}
Data & 2,903 days, Jan 2015 to Nov 2026\\
Convergence slope & $-0.4986$ (theory $-1/2$)\\
Dispersion ratio & $0.991$ to $1.021$\\
Ljung--Box $Q(10)$ & $1752$ squared vs $23.2$ crit.\\
Leverage corr & $-0.200$\\
MZ slope $b$ & $0.664$, $t_{b=1}=-3.70$\\
\textbf{It\^o vs actual P\&L} & $\mathbf{corr=+0.949}$\\
Backtest & $+36.06$ pts/cycle, $t=4.80$\\
Jump test & $Z=+1.617$, below $1.96$\\
\end{tabular}
\end{kb}

\begin{qb}[Q: ``Why use a model from 1973?'']
\scriptsize
The model is a \textbf{measurement framework}, not a claim about reality. The P\&L identity
comes from self-financing accounting plus It\^o plus the PDE. The $\sigma_t$ in it is the
\textbf{realised} volatility of whatever path actually happened. It never uses normality,
independence, or constant $\sigma$. We \emph{measured} the three ways the NIFTY breaks the
model (clustering, asymmetry, fat tails) and the identity still holds at $0.949$ on that same
data. \textbf{A model can be wrong about the distribution of returns and still be exactly right
about the accounting.}
\end{qb}

\columnbreak

\begin{qb}[Q: ``Which numerical methods did you use?'']
\scriptsize
Exact log-solution of GBM (so zero discretisation error) \textbullet\ Monte Carlo with a fixed
seed \textbullet\ closed-form Black--Scholes price, delta, gamma \textbullet\ central finite
differences, to verify the Greeks only \textbullet\ trapezoidal quadrature for the It\^o
integral \textbullet\ left-endpoint Riemann sums, as the It\^o convention requires \textbullet\
ordinary least squares written out by hand.

\textbf{If asked about Milstein:} in log coordinates $\sigma(t)$ is state-independent, so
$b'=0$ and the correction term vanishes. Milstein reduces to Euler here.
\end{qb}

\begin{tb}[Q: ``Did you use neural networks / PINNs?'']
\scriptsize
\textbf{No. None. Not anywhere.} There is no \texttt{torch}, \texttt{tensorflow},
\texttt{keras} or \texttt{sklearn} in the repository. This project is closed-form mathematics
plus Monte Carlo. \textbf{Nobody should hesitate on this question.}
\end{tb}

\begin{qb}[Q: ``Where is this used in the real world?'']
\scriptsize
Variance swaps are priced directly off this identity \textbullet\ the VIX is itself a
variance-swap calculation, which is why India VIX is the right thing to regress against
\textbullet\ an options market-maker's daily gamma P\&L is explained by this exact
decomposition \textbullet\ dispersion trading and vol-targeting funds harvest the same premium
\textbullet\ Carr and Wu (2009) document the variance risk premium across markets.
\end{qb}

\begin{qb}[Q: ``Is the data real?'']
\scriptsize
Yes. NIFTY~50 and India VIX daily closes, spot-checked against known history: the COVID bottom
of $7{,}610$ on 23 Mar 2020 and the India VIX spike to $84.6$ are both in the file. The
\textbf{simulated} paths appear only in the verification stage, deliberately, because testing
an estimator requires a case where the true answer is known in advance.
\end{qb}

\begin{tb}[Three things never to say]
\scriptsize
\textbf{1. Never claim an optimal rebalance frequency.} At zero cost the means are not
statistically distinguishable (spread $22.5$ against s.e.\ $11.9$). Frequency buys lower
dispersion, not a higher mean.

\textbf{2. Never say the premiums are market quotes.} They are Black--Scholes evaluated at
India VIX. Volunteer this before he asks.

\textbf{3. Never guess.} ``That part was done by X, let them answer'' costs far less than a
confident wrong answer.
\end{tb}

\begin{kb}[Volunteer these weaknesses first]
\scriptsize
The option premium is reconstructed, not quoted, so the P\&L \emph{level} inherits the model
(the measurement results do not) \textbullet\ costs are modelled, not measured \textbullet\
daily closes only, so the jump test's asymptotics are weak \textbullet\ one index,
at-the-money calls, flat $r$. The natural next version is the NSE option chain with its
bid--ask.
\end{kb}

\end{multicols}
\end{document}
